% Ерөнхий дүгнэлт

\phantomsection
\addchaptertocentry{Ерөнхий дүгнэлт}
\label{Summary}
\section*{Ерөнхий дүгнэлт}
Энэхүү дипломын ажлын хүрээнд байгууллагын 5S, бүтээмжийн хэрэгслүүд болон ажлын удирдлагыг нэгтгэсэн вэб, гар утасны платформыг зохион бүтээв. Үүнд:
\begin{itemize}
    \item lean, 5S, Gemba, Kaizen, шатлалт өдөр тутмын хурлын онолыг судалж, хэрэгслийг дадал болгоход сануулга, хяналтын механизм шаардлагатайг тодорхойлсон;
    \item ижил төстэй таван системийг харьцуулж, аудит, ажлын удирдлага, lean дадлыг нэг дор, монгол хэл дээр нэгтгэсэн шийдэл байхгүйг тогтоосон;
    \item бүтээмжийг хэмжих долоон үзүүлэлт ба ``өмнө--дараа'' үнэлгээний загварыг тодорхойлсон;
    \item 22 функциональ, 10 функциональ бус шаардлага, есөн use-case-ийн тодорхойлолт, шаардлагын мөрдөх матриц, класс, дараалал, үйл ажиллагааны диаграмм боловсруулсан;
    \item давхаргат архитектур, байршуулалтын диаграмм, 25 хүснэгттэй өгөгдлийн сан, дүрд суурилсан эрхийн загвар, зургаан гол алгоритмыг зохиомжилсон;
    \item backend, вэб, гар утасны аппыг хөгжүүлж, 2\,000 гаруй автомат тестээр баталгаажуулсан;
    \item таамаглалд суурилсан урьдчилсан эдийн засгийн тооцоо хийж, бодит үр ашгийг туршилтын ажиллагаагаар шалгахаар төлөвлөсөн.
\end{itemize}

Ажлын одоогийн төлөвийг гурван түвшинд ялгаж харуулав (Хүснэгт~\ref{table:status}). Системийн гол урсгалууд кодонд хэрэгжиж, автомат тестээр баталгаажсан боловч бодит байгууллагад ашиглаж бүтээмжийн өөрчлөлтийг хэмжих ажил хараахан хийгдээгүй; энэ нь дараагийн үзлэгүүдийн гол ажил юм.

\begin{table}[H]
\centering
\caption{Ажлын гүйцэтгэлийн төлөв}
\label{table:status}
\begin{tabular}{|p{5.6cm}|>{\centering\arraybackslash}p{2.6cm}|>{\centering\arraybackslash}p{2.6cm}|>{\centering\arraybackslash}p{2.6cm}|}
\hline
\textbf{Хэсэг} & \textbf{Кодонд хэрэгжсэн} & \textbf{Тестээр шалгасан} & \textbf{Туршилтаар хэмжсэн} \\ \hline
Ажил, төсөл, мэдэгдэл & \yes & \yes & -- \\ \hline
5S талбай, аудит, засах ажил & \yes & \yes & -- \\ \hline
Өдрийн ба долоо хоногийн тайлан & \yes & \yes & -- \\ \hline
Сарын тайлан, хаалт, архив & \yes & \yes & -- \\ \hline
Gemba, санаа, өглөөний хурал & \yes & \yes & -- \\ \hline
Сүлжээгүй горим (бичих) & $\sim$ & \yes & -- \\ \hline
KPI самбар (K1--K6) & $\sim$ & $\sim$ & -- \\ \hline
Хэрэглэгчийн судалгаа, SUS & -- & -- & -- \\ \hline
\end{tabular}
\end{table}

% TODO (Үзлэг 3): туршилтын ажиллагааны үр дүн, SUS оноо, цаашдын хөгжүүлэлтийг нэмнэ.
Цаашид системийг серверт байршуулж, туршилтын байгууллагад ажиллуулан хэрэглэхэд хялбар байдлыг SUS аргаар үнэлж, бодит хэмнэлтийг хэмжих, push мэдэгдэл нэмэх, апп дэлгүүрт нийтлэх ажлыг гүйцэтгэнэ.
